\chapter{Operational Label Definitions}\label{appendix-label-definitions}

This appendix consolidates the operational definitions used for annotation and
model evaluation. UI evaluability concerns what screenshots could establish in
principle. The verification label concerns what the recorded flow establishes
for a particular requirement.

\section{UI Evaluability}

\vspace{0.5\baselineskip}

\begin{longtable}{@{}>{\raggedright\arraybackslash}p{0.29\textwidth}>{\raggedright\arraybackslash}p{0.65\textwidth}@{}}
  \toprule
  \textbf{Label} & \textbf{Operational definition} \\
  \midrule
  \endfirsthead
  \toprule
  \textbf{Label} & \textbf{Operational definition} \\
  \midrule
  \endhead
  \texttt{UI\_VERIFIABLE} & All material claims can be assessed through rendered content or visible state changes. \\
  \texttt{PARTIALLY\_\allowbreak UI\_\allowbreak VERIFIABLE} & A material visible core exists, but full satisfaction also depends on hidden state, persistence, policy, business logic, or an external system. \\
  \texttt{NOT\_UI\_\allowbreak VERIFIABLE} & The central property has no stable visual manifestation or is too abstract for screenshot-based assessment. \\
  \bottomrule
\end{longtable}

\section{Verification Labels}

\vspace{0.5\baselineskip}

\begin{longtable}{@{}>{\raggedright\arraybackslash}p{0.29\textwidth}>{\raggedright\arraybackslash}p{0.65\textwidth}@{}}
  \toprule
  \textbf{Label} & \textbf{Use when} \\
  \midrule
  \endfirsthead
  \toprule
  \textbf{Label} & \textbf{Use when} \\
  \midrule
  \endhead
  \texttt{FULFILLED} & Every UI-observable core claim is visibly supported, at least one evidence unit is recorded, no observable core claim is contradicted, and no material uncertainty remains about visible behavior. \\
  \texttt{PARTIALLY\_\allowbreak FULFILLED} & At least one important claim is supported and at least one important claim remains missing, hidden, or ambiguous. A contradicted core claim instead requires \texttt{NOT\_FULFILLED}. \\
  \texttt{NOT\_FULFILLED} & Visible counter-evidence contradicts at least one observable core claim. Missing evidence alone is insufficient. \\
  \texttt{ABSTAIN} & The recorded screenshots support no reliable positive, partial, or negative decision. This is a completed semantic prediction, not a missing output. \\
  \bottomrule
\end{longtable}

Routine hidden dependencies do not block \texttt{FULFILLED} when the requested
visible outcome is shown. By contrast, an unverified system outcome or a
nontrivial hidden property blocks a full positive verdict.

\section{Uncertainty Reasons}

\vspace{0.5\baselineskip}

\begin{longtable}{@{}>{\raggedright\arraybackslash}p{0.39\textwidth}>{\raggedright\arraybackslash}p{0.55\textwidth}@{}}
  \toprule
  \textbf{Reason} & \textbf{Meaning} \\
  \midrule
  \endfirsthead
  \toprule
  \textbf{Reason} & \textbf{Meaning} \\
  \midrule
  \endhead
  \texttt{TEXTUAL\_AMBIGUITY} & The requirement wording is unclear or vague. \\
  \texttt{SCOPE\_OR\_CONTEXT\_\allowbreak AMBIGUITY} & The relevant screen, item, role, or flow segment is unclear. \\
  \texttt{QUANTIFIER\_OR\_\allowbreak COMPLETENESS\_\allowbreak AMBIGUITY} & Terms such as ``all,'' ``each,'' ``only,'' ``always,'' or ``closest'' are not visibly settled. \\
  \texttt{EVIDENCE\_\allowbreak INTERPRETATION\_\allowbreak AMBIGUITY} & Evidence is visible, but its meaning or causal interpretation is unclear. \\
  \texttt{FLOW\_COVERAGE\_GAP} & A necessary before, during, or after state is absent from the recorded flow. \\
  \texttt{UNVERIFIED\_SYSTEM\_\allowbreak OUTCOME} & A mechanism is shown, but the required success outcome is not. \\
  \texttt{NONTRIVIAL\_HIDDEN\_\allowbreak PROPERTY} & The requirement depends on a hidden property that screenshots cannot establish. \\
  \bottomrule
\end{longtable}

\section{Consistency Gates}

The following gates make the label policy testable:

\begin{itemize}
  \item \texttt{FULFILLED} requires recorded evidence and support for all
    observable core claims, and forbids a contradicted core claim or unresolved
    material uncertainty.
  \item \texttt{PARTIALLY\_FULFILLED} requires both meaningful support and an
    important missing, hidden, or ambiguous part.
  \item \texttt{NOT\_FULFILLED} requires visible counter-evidence against an
    observable core claim.
  \item \texttt{ABSTAIN} requires an insufficiency reason but no positive
    evidence.
\end{itemize}

At claim level, \texttt{SUPPORTED}, \texttt{CONTRADICTED}, \texttt{MISSING},
\texttt{HIDDEN}, \texttt{AMBIGUOUS}, and \texttt{OUT\_OF\_SCOPE} distinguish
positive evidence, counter-evidence, absent evidence, non-visible obligations,
unstable interpretations, and routine internal effects outside the visible
contract.
